\documentclass[10pt,a4paper]{amsart}
\usepackage[margin=19mm]{geometry}
\usepackage{amsmath,amssymb,mathtools}
\usepackage[scr=rsfs,cal=euler]{mathalfa}
\usepackage{xfrac}
\usepackage[unicode,hidelinks,bookmarks=false]{hyperref}
\newcommand{\ie}{i.e.\ }
\newcommand{\cf}{cf.\ }
\newcommand{\resp}{respectively\ }
\newcommand{\Subsection}{\subsection}
\providecommand{\nobreakdash}{\nobreak\hbox{-}}
\setlength{\emergencystretch}{2em}
\multlinegap0pt
\usepackage{amssymb}
\usepackage{enumerate}
\newtheorem{lem}{Lemma}
\title{Nilpotent groups have polynomially bounded homological filling invariants}
\author{Gabriel Pallier}
\date{Source: arXiv:2603.25890v1 (CC BY 4.0)}
\begin{document}
\maketitle

\noindent\emph{Source and licence.} Nilpotent groups have polynomially bounded homological filling invariants. Version arXiv:2603.25890v1 (CC BY 4.0), distributed by arXiv under the Creative Commons Attribution 4.0 International licence, http://creativecommons.org/licenses/by/4.0/, as recorded in the arXiv licence metadata for this exact version and shown on the abstract page https://arxiv.org/abs/2603.25890v1. Full licence notice: \texttt{licenses/arxiv-2603.25890-CC-BY-4.0.txt}.\par\smallskip

\noindent\textit{Editorial scope.} Counted unit: the article's lem lem:bounding-simple-vectors (source0.tex lines 105--113). The article's complete original proof of it is delivered with it (source0.tex lines 115--131, one \texttt{proof} environment of 17 lines). No mathematics is written, repaired or re-derived: the statement and the proof are the article's own lines, in the article's own notation, and no retained line is substituted or re-flowed.  No further source-local context is needed: every cross-reference of every retained line resolves inside the two delivered spans.  Nothing was omitted from any retained span, so the package carries no omission record.  The retained lines cite 1 works, whose entries are transcribed verbatim from the article's own bibliography source in the e-print and delivered with short editorial labels recorded in the item's reference mapping.  No retained line contains a figure environment, an image-inclusion command or drawing code, so no image is delivered and the package declares no asset dependency.  Editorial formatting only, in addition: an AMS \texttt{amsart} preamble that restates the article's own declarations and macros used by the retained lines, namely the environments \texttt{lem} on separate counters, together with the standard packages those lines need.  The retained environments print in this document as Lemma 1 in order of appearance, whereas the article prints the same items with the numbering of its own full document.  The complete native source of the article as distributed in the arXiv e-print for this exact version is bundled unchanged as \texttt{sources/197249/source0.tex}.
\begin{lem}[Polynomial similarity]
\label{lem:bounding-simple-vectors}
With notation as above, there exists $R_d\in \mathbf R[t]$, nondecreasing on $[0,+\infty)$ and such that for any $d$-vector field $V = V_1 \wedge \cdots \wedge V_d$ on $\mathbf R^d$ we have, for all $x = (x_1,\ldots, x_n) \in \mathbf R^n \simeq \mathfrak g$,
\begin{equation}
\label{eq:bounding-simple-vectors}
\frac{\Vert V_{x} \Vert_{\Lambda^dh_0}}{R_d (\Vert x \Vert)}
    \leqslant \Vert  V_{x} \Vert_{\Lambda^d \exp^\star h} \leqslant R_d (\Vert x \Vert) \Vert  V_{x} \Vert_{\Lambda^d h_0}.
\end{equation}    
\end{lem}

\begin{proof}

The Baker-Campbell-Hausdorff series for $G$ only has finitely many non-vanishing terms because $G$ is nilpotent, and so the left-invariant vector fields $X_1,\ldots, X_n$ which coincide with $e_1, \ldots,e_n$ at the origin are pulled-back to the Lie algebra via the exponential as
\begin{equation}
\label{eq:bounding-simple-vectors-1}
     (\exp_G^\star X_i)(x_1,\ldots, x_n) = e_i + \sum_{j>i} P_{i,j}(x_1,\ldots, x_n) e_j 
\end{equation}
where $P_{i,j} \in \mathbf R[x_1,\ldots, x_n]$ are polynomials of total degrees at most $s -1$ (see e.g. \cite{LeDonneTripaldiCCLow} for many concrete examples of such calculations). 
Fix now an integer $d \in \{1,\ldots, n\}$. We will generalize the previous considerations to the $d$-areas. Denote by $I = \{i_1, \ldots, i_d \}$ a multi-index with $1 \leqslant i_1 < \cdots < i_d \leqslant n$ and equip the set of multiindices with the lexicographic order.  An orthonormal frame is given by the simple left-invariant $d$-vectors $X_I = X_{i_1} \wedge \cdots \wedge X_{i_d}$ and
\begin{equation}
     \operatorname{exp}^\star X_I = e_I + \sum_{J > I} P^{[d]}_{I,J}(x_1,\ldots, x_n) e_J, \notag
\end{equation}
where $P_{I,J}^{[d]}$ are again polynomials.
 This equation expresses that the transition matrix between two orthogonal frames of $\Lambda^d h_0$ and 
$\Lambda^d \exp^\star h$ is unipotent with its non-diagonal entries being polynomials in $(x_1, \ldots, x_n)$. Inverting this matrix, we check that its inverse has the same properties.
The bounds \eqref{eq:bounding-simple-vectors} follow from this observation for some (maybe not nondecreasing) polynomial $\widetilde R_d$ instead of $R_d$; we may finally replace $\widetilde R_d$ by a larger polynomial $R_d$ which is nondecreasing on $[0,+\infty)$ in order to reach the conclusion of the lemma.
\end{proof}

\begin{thebibliography}{99}
\bibitem[LeDonn]{LeDonneTripaldiCCLow} Enrico Le~Donne and Francesca Tripaldi. \newblock A cornucopia of {C}arnot groups in low dimensions. \newblock {\em Anal. Geom. Metr. Spaces}, 10(1):155--289, 2022.
\end{thebibliography}
\end{document}
